\section*{第 \ref{ch:b3:crystalline-solids} 章——晶体固体}

\begin{solution}{exo:b3:crystalline-solids:1}
(a) 简单立方：$8\times\tfrac18 = 1$；体心立方：$8\times\tfrac18 + 1 = 2$；
面心立方：$8\times\tfrac18 + 6\times\tfrac12 = 4$。(b) 体对角线
$a\sqrt3 = 4r$：$r = a\sqrt3/4$。(c) 面对角线 $a\sqrt2 = 4r$：
$r = a\sqrt2/4$。(d) $0.68$ 与 $0.74$：
\omterm{rem:b3:crystalline-solids:bonds}{金属键}无方向性，要求最大堆积，因此铜、铝、银、金采用面心立方（或同样致密的
六方）结构。
\end{solution}

\begin{solution}{exo:b3:crystalline-solids:2}
(a) $a/\sqrt2 = \qty{2.55}{\angstrom}$。(b) 12——密堆积的
配位数。(c) \qty{8960}{kg/m^3}
（\cref{ex:b3:crystalline-solids:density}）。(d) $n = 4/a^3 =
\qty{8.5e28}{m^{-3}}$，因此每立方微米中含
$8.5\times10^{10}$ 个原子——这正是芯片金属工艺
靠统计、而非靠木工的原因。
\end{solution}

\begin{solution}{exo:b3:crystalline-solids:3}
(a) $\sin\theta_1 = \lambda/2d = 0.273$：$\theta_1 =
15.8^\circ$。(b) $n \le 2d/\lambda = 3.66$：有三个级次。(c)
$\sin\theta \le 1$ 要求 $n\lambda \le 2d$；可见光的
$\lambda \sim \qty{5000}{\angstrom}$ 比晶面间距长一千倍，
——任何晶面间距都无法使其衍射。(d) $\sin\theta$
下降 1\,\%：$\theta_1 \to 15.7^\circ$。衍射仪是一台精密的
温度计——这一微小偏移正是在原子尺度上测量热膨胀的
方法。
\end{solution}

\begin{solution}{exo:b3:crystalline-solids:4}
(a) $d_{100} = a > d_{110} = a/\sqrt2 > d_{111} = a/\sqrt3$。(b)
$d$ 最大者，角度最小：(100)。(c) 八面体
(111) 族——在金刚石结构中，其原子层两两配对，形成
键合很强的双层，双层之间是宽阔而键合稀疏的间隙：
这就是解理面。(d) 每个晶粒在相同的 $\theta$ 下反射，
但方位角随机：反射光线散成半顶角为
$2\theta$ 的锥面，在探测器上切出环纹。
\end{solution}

\begin{solution}{exo:b3:crystalline-solids:5}
(a) 干涉需要光程差与 $\lambda$ 同量级；
间距为埃量级，故 $\lambda$ 也必须如此。(b) $E =
hc/\lambda \approx \qty{8}{keV}$。(c) $p = h/\lambda$，$V =
p^2/2m_{\text{e}}e \approx \qty{60}{kV}$（相对论效应使其
减小百分之几）——这正是电子显微镜的工作电压。(d)
$\lambda = h/\sqrt{3m_{\text{n}}k_{\text{B}}T} \approx
\qty{1.5}{\angstrom}$：室温中子天生便适于衍射。
它们与原子核而非电子云发生散射，
——因此能清楚地看到氢原子，并带有磁矩，
可用于绘制磁有序；这两点对 X 射线几乎都不可见。
\end{solution}

\begin{solution}{exo:b3:crystalline-solids:6}
(a) 每个体积为 $a^3$ 的\omterm{def:b3:crystalline-solids:lattice}{晶胞}含四对原子。(b) $N_{\text{A}} =
4M/\rho a^3 = 4\times0.05844/(2165\times\qty{1.794e-28}{})
\approx \qty{6.02e23}{mol^{-1}}$。(c) $N_{\text{A}} \propto
a^{-3}$：$a$ 的 0.1\,\% 误差导致 $N_{\text{A}}$ 的 0.3\,\% 误差。(d)
该方法通过假定每个\omterm{def:b3:crystalline-solids:lattice}{晶胞}都充满并
且全同来计数原子：空位、杂质与镶嵌边界都会造成计数偏差
——因此千克重新定义中那些近乎苛求地完美的
硅球才如此讲究。
\end{solution}

\begin{solution}{exo:b3:crystalline-solids:7}
(a) 每个离子所感受到的为 $2\sum_{n\ge1}(-1)^{n+1}/n = 2\ln2$，以
$e^2/4\pi\varepsilon_0r_0$ 为单位（因子 2：两侧贡献），为吸引作用。(b)
$e^2/4\pi\varepsilon_0r_0 = \qty{5.11}{eV}$；$\times1.7476 =
\qty{8.93}{eV}$。(c) 最小化时消去 $1/9$ 的吸引：
$U = 8.93\times\tfrac89 = \qty{7.94}{eV}$。(d) 与
\qty{7.9}{eV} 相差不到百分之一：点电荷晶格加上一个刚度指数，
即可解释离子\omterm{def:b3:crystalline-solids:lattice}{晶体}的全部能量预算
——量子力学隐藏在 $r_0$ 与该指数之中。
\end{solution}

\begin{solution}{exo:b3:crystalline-solids:8}
(a) $r_0 = 2^{1/6}\times3.40 = \qty{3.82}{\angstrom}$，比测得的
\qty{3.76}{\angstrom} 大 1.5\,\%（每个原子还受到其十二个近邻的作用，
使势阱变深）。(b)
$T_{\text{m}} \sim \epsilon/2k_{\text{B}} \approx \qty{60}{K}$：
恰为 \qty{84}{K} 的合适量级。(c) 氦的\omterm{thm:b3:harmonic-oscillator:spectrum}{零点能}
$\hbar\omega/2$ 在如此浅的势阱中与势阱本身相当：
\omterm{def:b3:crystalline-solids:lattice}{晶体}自行瓦解，因此氦在 $T = 0$ 时仍为液体，
除非加压。(d) 势阱深度 $\sim\qty{0.01}{eV}$，而\omterm{rem:b3:crystalline-solids:bonds}{共价键}约为
$\sim\qty{3.6}{eV}$ 每键：能量相差三百倍，正是
从霜冻到金刚石的全部距离。
\end{solution}

\begin{solution}{exo:b3:crystalline-solids:9}
(a) $ka \ll 1$：$\omega \approx 2\sqrt{K/m}\,(ka/2) =
a\sqrt{K/m}\,k$。(b) 将\omterm{def:b3:crystalline-solids:lattice}{晶胞}拉伸 $\delta$：应力 $\sim
K\delta/a^2$，应变 $\delta/a$，故 $E \sim K/a$，即 $K \sim
Ea$。(c) $K \approx \qty{31}{N/m}$，$v = a\sqrt{K/m} \approx
\qty{4400}{m/s}$——与测得的 \qty{4000}{m/s} 相差在 10\,\% 以内，
其弹簧常数由\omterm{thm:b3:continuum-elasticity:hooke}{杨氏模量}估算而得。(d) $v = a\sqrt{K/m}$：
键越硬、原子越轻，速率越大——金刚石两者都取到最大，因此
声速达 \qty{18000}{m/s}，$\Theta_{\text{D}} \approx
\qty{2200}{K}$。
\end{solution}

\begin{solution}{exo:b3:crystalline-solids:10}
(a) $\dd\omega/\dd k \propto \cos(ka/2) = 0$，在 $k = \pi/a$ 处。(b)
$u_n \propto (-1)^n$：每个原子都与两侧近邻反相——
形成驻波，能量在原位来回振荡。(c)
当 $\lambda = 2a$ 时，相继原子反向散射的波，
光程差恰为一个波长：这正是
沿链本身的布拉格条件。晶格反射了它自身的振动，
前行与后行的波锁定为驻波图样，
行波因此无法继续传播。(d) $\omega_{\text{最大}} = 2\sqrt{K/m} \approx
\qty{3.4e13}{rad/s}$，即 $\sim\qty{5}{THz}$：远红外
——晶格振动与红外光在同一个倍频程内相遇，
这正是下一题所依据的事实。
\end{solution}

\begin{solution}{exo:b3:crystalline-solids:11}
（a）在光学支中，$+$ 与 $-$ 两个子晶格沿相反
方向运动：一个振荡的电偶极子，正是光波
所作用的对象。（b）$\omega \approx \sqrt{2K/\mu}$，其中 $\mu =
\qty{2.3e-26}{kg}$：$\omega \approx \qty{4.7e13}{rad/s}$，
$\lambda = 2\pi c/\omega \approx \qty{40}{\micro m}$——属于远
红外波段（食盐实测的剩余射线带位于
$\qty{61}{\micro m}$，双弹簧模型给出的量级正确）。（c）
在该波段，\omterm{def:b3:crystalline-solids:lattice}{晶体}自身的共振吸收并再次反射
该波：透明的食盐转为镜面式不透明。（d）每个
\omterm{def:b3:crystalline-solids:lattice}{晶胞}只含一个原子，意味着没有第二个子晶格与之相互拍频：仅有单一支，
因而没有能隙。
\end{solution}

\begin{solution}{exo:b3:crystalline-solids:12}
（a）$N$ 个原子 $\times$ 3 个位移方向 = $3N$
个振子，因此线性谱在计入
$3N$ 个模式后截止：由此定义 $k_{\text{D}}$。（b）
$k_{\text{D}} = (6\pi^2n)^{1/3} = \qty{1.7e10}{m^{-1}}$，
$\Theta_{\text{D}} = \hbar vk_{\text{D}}/k_{\text{B}} \approx
\qty{350}{K}$——与由量热测得的 \qty{343}{K} 相符，
由声速和密度得出。（c）直线遗漏了布里渊区边界处的平坦化：
\omterm{thm:b3:photons-phonons:debye}{德拜模型}高估了高频模式，因此拟合
在中间温度下有偏差；$T^3$ 定律（仅适用于长
波）是可靠的。（d）两者都表明，谱在波长
达到原子间距时结束——每个原子对应一个模式，
无论表现为截止波矢还是截止温度。
\end{solution}

\begin{solution}{pb:b3:crystalline-solids:1}
\textbf{1.} $2d\sin\theta = n\lambda$。单色：一个
$\lambda$ 使每个环角与唯一的面间距对应。粉末：在
随机取向的晶粒中总有一些恰好满足反射条件——每一
族晶面都同时给出信号。
\textbf{2.} $\theta = 19.05^\circ, 22.15^\circ, 32.2^\circ,
38.75^\circ$；$\sin^2\theta = 0.107, 0.142, 0.284, 0.392$。
\textbf{3.} 将 $d = a/\sqrt{h^2+k^2+l^2}$ 代入布拉格
公式（$n = 1$）：$\sin^2\theta = (\lambda^2/4a^2)(h^2+k^2+l^2)$。
\textbf{4.} 比值 $1 : 1.33 : 2.67 : 3.68$；乘以 3：
$3.0 : 4.0 : 8.0 : 11.0$。
\textbf{5.} 全奇数 $(111)$ 以及全偶数 $(200)$、$(220)$、
$(311)$ 恰好给出 3、4、8、11；缺失的 1、2、5、6、7 号峰
排除了简单立方，其图样本应显示它们。
\textbf{6.} 体心立方允许的偶数和给出比值
$1:2:3:4$——节律不同。比值在不知道
$a$、波长漂移和样品取向偏差时依然成立：奇偶性是几何性质，
而非标定的结果。
\textbf{7.} 晶面数少时，干涉极大较钝，类似
只有五条缝的光栅：环宽 $\sim$ 微晶
尺寸的倒数（谢乐公式）——锐利的环说明晶粒
生长良好。
\textbf{8.} $a = \lambda\sqrt{h^2+k^2+l^2}/2\sin\theta = 4.088,
4.086, 4.089, 4.082\,\unit{\angstrom}$。
\textbf{9.} $a \approx \qty{4.09}{\angstrom}$。
\textbf{10.} $\rho = 4M/N_{\text{A}}a^3$。
\textbf{11.} $4/N_{\text{A}}a^3 = \qty{9.74e4}{mol/m^3}$：
银 $\to \qty{10.5}{g/cm^3}$（与银的手册值相符），
金 $\to 19.2$（与金的手册值相符！），铂 $\to 19.0$
（与铂的 21.4 矛盾——其真实的 $a$ 为
\qty{3.92}{\angstrom}）。铂被排除；银和金都保留下来，
因为它们的\omterm{def:b3:crystalline-solids:lattice}{晶格常数}几乎完全相同。
\textbf{12.} 需要基于质量的测量：密度（或仅摩尔
质量）。银和金的 $a$ 相差不超过 0.2\,\%，但它们的原子
质量相差 1.8 倍——任何晶格上的巧合都无法
弥补这一差异。
\textbf{13.} 依阿基米德原理测量金锭：密度接近 \qty{10.5}{g/cm^3}
$\to$ 银，与其锈蚀情况相符（银会变暗；
金则会保持光亮）。微观上的\omterm{def:b3:crystalline-solids:lattice}{晶胞}质量与
宏观称量结果一致——同一个
$\rho = 4M/N_{\text{A}}a^3$ 从两个方向读取。
\textbf{14.} 原子在\emph{平均
位置}附近振动：决定环角度的平均晶格保持不变；
随机位移只会重新分配强度。
\textbf{15.} 强度随 $T$ 平滑下降（
$\exp(-q^2\langle x^2\rangle)$ 型德拜–沃勒因子），
损失的光以环与环之间的漫射背景重新出现。
\textbf{16.} $K \approx Ea = \qty{8.3e10}{}\times
\qty{2.89e-10}{} \approx \qty{24}{N/m}$。
\textbf{17.} $v = a\sqrt{K/m} \approx \qty{3300}{m/s}$ ---
与实测的 \qty{2700}{m/s} 相差 20\,\%。
\textbf{18.} $\langle x^2\rangle = 3k_{\text{B}}T/K$：均方根
$\approx\qty{0.23}{\angstrom}$，约占键长的 8\,\%
在室温下已是如此。
\textbf{19.} 10\,\% 判据外推得到 $\sim\qty{500}{K}$
而真实值为 \qty{1235}{K}：量级正确，偏差约两倍——
——我们的单弹簧 $K$ 偏软，林德曼判据是一条标度
规则，而非定律。结论依然成立：熔化发生于热振动
可与晶格本身相抗衡之时。
\textbf{20.} $d_{110} = a/\sqrt2 = \qty{2.04}{\angstrom}$：
$\sin\theta = 0.378$，$2\theta \approx 44.5^\circ$。
\textbf{21.} 耦合强度：电子（带电）在
纳米尺度上散射——适用于表面与薄膜；X 射线在微米——
尺度上散射——适用于块体粉末；中子在厘米尺度上——适用于整个发动机部件。
\textbf{22.} 环收缩为一个或两个宽阔的光晕：
液体保留短程有序（偏好的近邻
距离），但不存在长程有序——没有剩余的周期结构可以
产生锐利的干涉。
\textbf{23.} DNA——照片 51。衍射可读出重复
距离，远低于任何光学显微镜的分辨极限；X 射线图样
暴露出螺旋结构，\qty{3.4}{\angstrom} 的模糊条纹对应其堆叠的
碱基对。
\textbf{24.} 结构是确定有序的（
准周期性——有序却不重复）：锐利的斑点需要
长程相位相干性，而非周期性——这一发现
拓宽了\omterm{def:b3:crystalline-solids:lattice}{晶体}学对\omterm{def:b3:crystalline-solids:lattice}{晶体}本身的定义。
\textbf{25.} 比值 $3:4:8:11 \to$ 面心立方；$a = \qty{4.09}{\angstrom}$
$+$ 密度 $\to$ 银，而非金或铂；由于加热会使环变淡，
熔化则会抹去环，故应在冷态下对金锭做 X 射线检测——
博物馆的银，经干涉认证。
\end{solution}
